% Formalized in GSQG/PatchEquation/*.lean (main theorem: GSQG/PatchEquation/MainTheorem.lean).
\label{sec:c13_patch_equation}

This chapter derives the graph equation from a two-dimensional weak patch system and proves
its converse for the smooth unbounded graphs used in \cref{thm:main}.
All declarations of this chapter are formalized in Lean.
The scope is weak patch solutions with smooth interfaces before the corner time $T$: no classical
vector velocity on the interface and no continuation beyond $T$ is asserted.
For background on contour dynamics and the divergent tangential trace see
\cite[Section~3.1]{Hunter-Shu:regularized-approximate-model-sqg}; their parameter equals $2-\alpha$
in the convention used here. The rotation, density orientation and time normalization are fixed explicitly below.

We work with $1<\alpha<2$ and with a positive kernel coefficient $c$ until the normalization in
\cref{sec:c13_operator}. This includes $\alpha=p+1\in(7/4,2)$ in the application.
The space $E=\R^2$ has its Euclidean norm and Lebesgue measure. In Lean, use
\texttt{EuclideanSpace R (Fin 2)} and explicit volume-preserving coordinate maps;
the default norm of a product of two real normed spaces is not the Euclidean radius.
Write $J(x_1,x_2)=(-x_2,x_1)$, so $\nabla^\perp=J\nabla$.
Spatial and spacetime test functions are smooth and compactly supported in their indicated open domains.
Constants in quantitative statements are chosen before the variables in which uniformity is asserted.

\section{The two-dimensional operator and its normalization}
\label{sec:c13_operator}

\begin{definition}[Kernel, adjoint cancellation and renormalized potential]
  \label{def:c13_kernel_data}
  \lean{GSQG.c13_E, GSQG.c13_pt, GSQG.c13g_aux_div, GSQG.c13_rotation, GSQG.c13_potential_kernel, GSQG.c13_velocity_kernel, GSQG.c13_adjoint, GSQG.c13_renormalized_potential}
  Let $E=\R^2$ with the Euclidean norm, and write $(x,y)\in E$ for the point with coordinates $x,y$.
  For a vector field $Q:E\to E$ put $\operatorname{div}Q=\partial_1Q_1+\partial_2Q_2$.
  For $z\ne0$ set
  \[
    P_{\alpha,c}(z)=c|z|^{-\alpha},\qquad
    K_{\alpha,c}(z)=\alpha c\,Jz\,|z|^{-\alpha-2}=-J\nabla P_{\alpha,c}(z),
  \]
  and set both kernels to zero at the single point zero.
  For a vector test function $\Phi:E\to E$ define
  \[
    B^*_{\alpha,c}\Phi(Y)=\int_E K_{\alpha,c}(z)\cdot(\Phi(Y+z)-\Phi(Y))\,dz.
  \]
  For a bounded measurable real density $\theta$ define
  \[
    V_{\alpha,c}[\theta](X)=\int_E
       [P_{\alpha,c}(X-Y)-\mathbf1_{|Y|>1}P_{\alpha,c}(-Y)]\theta(Y)\,dY.
  \]
  These are total formulas; the next lemmas establish their analytic meanings.
  The unsubtracted integral $P*\theta$ is not asserted to converge for an infinite patch.
\end{definition}

\begin{definition}[Canonical and front normalizations]
  \label{def:c13_normalization}
  \lean{GSQG.c13_riesz_factor, GSQG.c13_standard_coefficient, GSQG.c13_time_ratio}
  \uses{def:frac_multiplier}
  Put
  \[
    A_\alpha=\frac{2^{2-\alpha}\pi\Gamma(1-\alpha/2)}{\Gamma(\alpha/2)},\qquad
    c_{\rm std}=A_\alpha^{-1},\qquad c_{\rm front}=\operatorname{calpha}(\alpha-1),\qquad
    \kappa_\alpha=\frac{c_{\rm front}}{c_{\rm std}}.
  \]
  The two-dimensional Fourier convention is
  $\widehat f(\xi)=\int_E f(X)e^{-iX\cdot\xi}dX$, with inverse factor $(2\pi)^{-2}$.
  The canonical operator $-\nabla^\perp(-\Delta)^{-1+\alpha/2}$ has multiplier
  $-iJ\xi|\xi|^{\alpha-2}$, given value zero at $\xi=0$.
  Identification of the real-space operator with this multiplier is a theorem below.
  The two coefficients are kept distinct; the final application records the corresponding time change.
\end{definition}

\begin{lemma}[Euclidean coordinates and radial majorants]
  \label{lem:c13_radial_majorants}
  \lean{GSQG.c13_radial_majorants}
  \uses{def:c13_kernel_data}
  Coordinate identification $E\leftrightarrow\R\times\R$ preserves Lebesgue measure, and
  $|z|^2=z_1^2+z_2^2$. For $r>0$,
  \[
    \int_{|z|<r}|z|^{-\alpha}dz=\frac{2\pi}{2-\alpha}r^{2-\alpha},\qquad
    \int_{|z|>r}|z|^{-\alpha-1}dz=\frac{2\pi}{\alpha-1}r^{1-\alpha}.
  \]
  The displayed nonnegative integrands are integrable on their respective domains.
  Consequently there is $C_{\alpha,c}>0$ such that for every $\xi\in E$,
  \[
    \int_E|K_{\alpha,c}(z)|\,|e^{-iz\cdot\xi}-1|dz
       \le C_{\alpha,c}|\xi|^{\alpha-1},
  \]
  with the integral zero when $\xi=0$.
\end{lemma}
\begin{proof}
  Use the volume-preserving Euclidean-to-coordinate map and radial integration in dimension two.
  For the last bound split at $r=|\xi|^{-1}$ and use $|e^{-is}-1|\le\min(|s|,2)$.
  The cases $\alpha>1$ and $\alpha<2$ control different ends of the integral.
\end{proof}

\begin{lemma}[$L^1$ translation estimate for test fields]
  \label{lem:c13_test_translation}
  \lean{GSQG.c13_test_translation}
  For every vector test function $\Phi:E\to E$ and every $z\in E$,
  \[
    \int_E|\Phi(Y+z)-\Phi(Y)|dY
       \le\min\left(|z|\int_E\|D\Phi(Y)\|dY,\ 2\int_E|\Phi(Y)|dY\right).
  \]
  All the displayed integrals are finite.
\end{lemma}
\begin{proof}
  Apply segment FTC to $\Phi(Y+tz)$, integrate its norm, and use Fubini and translation invariance.
  The second bound is the triangle inequality. This is a finite-dimensional test-function lemma.
\end{proof}

\begin{lemma}[Absolute adjoint pairing]
  \label{lem:c13_adjoint_integrable}
  \lean{GSQG.c13_adjoint_integrable}
  \uses{def:c13_kernel_data}
  For every vector test $\Phi$, the function
  $(Y,z)\mapsto K_{\alpha,c}(z)\cdot(\Phi(Y+z)-\Phi(Y))$ is absolutely integrable on $E\times E$.
  The inner integral defining $B^*\Phi(Y)$ converges absolutely for every $Y$, and
  \[
    \|B^*\Phi\|_{L^1(E)}\le C_{\alpha,c}
       (\|\Phi\|_{L^1(E)}+\|D\Phi\|_{L^1(E)}).
  \]
  Thus $\int\theta B^*\Phi$ is absolutely convergent for every bounded measurable $\theta$.
\end{lemma}
\begin{proof}
  \uses{lem:c13_radial_majorants, lem:c13_test_translation}
  Use Tonelli for the norm majorant and split $|z|\le1$ from its complement.
  Pointwise convergence uses the bounded derivative of the test field near zero and its bounded values at infinity.
\end{proof}

\begin{lemma}[Test potential and adjoint divergence identity]
  \label{lem:c13_test_potential}
  \lean{GSQG.c13_test_potential}
  \uses{def:c13_kernel_data}
  For a vector test $\Phi$, let $A_\Phi(Y)=\int_E P_{\alpha,c}(X-Y)\Phi(X)dX$.
  This integral is absolutely convergent, $A_\Phi$ is $C^1$, and
  \[
    B^*\Phi=-\operatorname{div}(J A_\Phi).
  \]
  If $\operatorname{supp}\Phi\subset\{|X|\le L\}$, then for $|Y|>2L+1$,
  $|A_\Phi(Y)|\le C_{\alpha,c}|Y|^{-\alpha}\|\Phi\|_1$.
  In particular $B^*(\nabla\zeta)=0$ for every scalar test $\zeta$.
\end{lemma}
\begin{proof}
  \uses{lem:c13_radial_majorants, lem:c13_adjoint_integrable}
  Differentiate the compact test factor after shifting variables. Integration by parts in the
  locally integrable potential kernel gives the adjoint identity; retain the cancellation at zero.
  For a gradient test, the scalar curl is zero by equality of mixed derivatives.
  Lean route: $A_\Phi$ is the convolution $P*\Phi$ ($P$ even), differentiated on the compact test factor;
  the adjoint identity $B^*\Phi(Y)=\int P(z)\operatorname{curl}\Phi(Y+z)dz$ is proved for the regularized kernel
  $P_\delta=c(|z|^2+\delta^2)^{-\alpha/2}$ by whole-space integration by parts and oddness of $K_\delta$,
  then $\delta\to0$ by dominated convergence, keeping the cancellation $\Phi(Y+z)-\Phi(Y)$ until the limit.
\end{proof}

\begin{lemma}[Absolute convergence of the renormalized potential]
  \label{lem:c13_potential_convergence}
  \lean{GSQG.c13_potential_convergence}
  \uses{def:c13_kernel_data}
  For every $R\ge1$ there is $C_{\alpha,c,R}>0$ such that for every bounded measurable $\theta$
  with $|\theta|\le M$, $M\ge0$, and every $|X|\le R$,
  \[
    \int_E|P(X-Y)-\mathbf1_{|Y|>1}P(-Y)|\,|\theta(Y)|dY\le C_{\alpha,c,R}M.
  \]
  The integral defining $V[\theta]$ is therefore absolutely convergent at every $X$, measurable
  and locally bounded. Changing the subtraction radius to another positive radius changes it
  by a finite constant independent of $X$.
\end{lemma}
\begin{proof}
  \uses{lem:c13_radial_majorants}
  On $|Y|>2R+2$, use the mean-value estimate
  $|P(X-Y)-P(-Y)|\le C_{\alpha,c}R|Y|^{-\alpha-1}$.
  In the remaining bounded region integrate the translated singular kernel using its radial
  integral. Measurability follows from parameter integration. The difference of two subtraction
  choices is an integral over a bounded annulus and contains no $X$.
\end{proof}

\begin{lemma}[Renormalized potential and its distributional gradient]
  \label{lem:c13_renormalized_potential}
  \lean{GSQG.c13_renormalized_potential_properties}
  \uses{def:c13_kernel_data}
  For every bounded measurable $\theta$, the defining integral for $V_{\alpha,c}[\theta](X)$
  converges absolutely for every $X$ and defines a locally bounded measurable function.
  For every vector test $\Phi$,
  \[
    \int_E\theta B^*\Phi=-\int_E V_{\alpha,c}[\theta]\,\operatorname{div}(J\Phi).
  \]
  Both sides converge absolutely. This is $B_{\alpha,c}\theta=-J\nabla V_{\alpha,c}[\theta]$
  in distributions. Replacing the subtraction radius $1$ by another positive radius changes
  $V$ by a finite constant independent of $X$ and leaves the velocity unchanged.
\end{lemma}
\begin{proof}
  \uses{lem:c13_potential_convergence, lem:c13_adjoint_integrable, lem:c13_test_potential}
  For $|Y|>2|X|+2$, the subtracted kernel is bounded by
  $C_{\alpha,c}|X||Y|^{-\alpha-1}$. Near $Y=X$ use local integrability of $|X-Y|^{-\alpha}$.
  Establish local bounds by integrating the singular kernel, not by taking a pointwise supremum over its moving pole.
  Integrate by parts in the kernel with a small-radius cutoff; its error is $O(r^{2-\alpha})$.
  Keep the subtraction until Fubini is justified. Its contribution against $\operatorname{div}(J\Phi)$
  then vanishes because the integral of a compactly supported divergence is zero.
\end{proof}

\begin{lemma}[Gamma layer representation of the singular kernel]
  \label{lem:c13_gamma_kernel_layers}
  \lean{GSQG.c13_gamma_kernel_layers}
  \uses{def:c13_kernel_data}
  For $z\ne0$,
  \[
    |z|^{-\alpha}=\frac1{\Gamma(\alpha/2)}\int_0^\infty
        s^{\alpha/2-1}e^{-s|z|^2}ds,\qquad
    K_{\alpha,c}(z)=\frac{2c}{\Gamma(\alpha/2)}Jz\int_0^\infty
        s^{\alpha/2}e^{-s|z|^2}ds.
  \]
  Both scalar parameter integrals converge absolutely. For each $\xi\in E$, multiplying the
  second representation by $e^{-iz\cdot\xi}-1$ yields an absolutely integrable double integrand
  in $(s,z)\in(0,\infty)\times E$.
\end{lemma}
\begin{proof}
  \uses{lem:c13_radial_majorants}
  Substitute $u=s|z|^2$ in the Gamma integral and use
  $\Gamma(1+\alpha/2)=(\alpha/2)\Gamma(\alpha/2)$.
  Tonelli on the norm reduces double-integrability to the radial cancellation majorant.
\end{proof}

\begin{lemma}[Two-dimensional Gaussian transform and first moment]
  \label{lem:c13_gaussian_first_moment}
  \lean{GSQG.c13_gaussian_first_moment}
  For every $s>0$ and $\xi\in E$, the following integrals converge absolutely:
  \[
    \int_Ee^{-s|z|^2-i z\cdot\xi}dz=\frac\pi s e^{-|\xi|^2/(4s)},\qquad
    \int_E z e^{-s|z|^2-i z\cdot\xi}dz
       =-\frac{i\pi}{2s^2}\xi e^{-|\xi|^2/(4s)},\qquad
    \int_E z e^{-s|z|^2}dz=0.
  \]
  The vector formulas are interpreted componentwise after complexification.
\end{lemma}
\begin{proof}
  Apply the inner-product-space Gaussian Fourier integral in dimension two with
  phase coefficient $-i$, or coordinate Fubini and the one-dimensional formula.
  Differentiate in $\xi$ under the integral, with a Gaussian moment majorant, for the first moment.
  Oddness gives the final identity. Keep the $e^{-iz\cdot\xi}$ convention explicit.
\end{proof}

\begin{lemma}[Inverse-parameter Gamma integral]
  \label{lem:c13_inverse_gamma_integral}
  \lean{GSQG.c13_inverse_gamma_integral}
  For every $a,b>0$, the integrand is integrable on $(0,\infty)$ and
  \[
    \int_0^\infty s^{-1-a}e^{-b/s}ds=b^{-a}\Gamma(a).
  \]
\end{lemma}
\begin{proof}
  Use the decreasing substitution $u=b/s$, including the absolute Jacobian.
  Gamma integrability proves convergence at both endpoints before the change of variables.
\end{proof}

\begin{lemma}[Gaussian evaluation of the velocity symbol]
  \label{lem:c13_kernel_symbol}
  \lean{GSQG.c13_kernel_symbol}
  \uses{def:c13_kernel_data, def:c13_normalization}
  For every $\xi\in E$, the vector integral is absolutely convergent and
  \[
    \int_E K_{\alpha,c}(z)(e^{-iz\cdot\xi}-1)dz
       =-icA_\alpha J\xi|\xi|^{\alpha-2},
  \]
  where the right side is zero at $\xi=0$. In particular $A_\alpha,c_{\rm std},\kappa_\alpha>0$.
\end{lemma}
\begin{proof}
  \uses{lem:c13_gamma_kernel_layers, lem:c13_gaussian_first_moment, lem:c13_inverse_gamma_integral}
  Use the Gamma representation
  $K(z)=\frac{2c}{\Gamma(\alpha/2)}Jz\int_0^\infty s^{\alpha/2}e^{-s|z|^2}ds$.
  Absolute Fubini follows from the radial cancellation majorant before evaluating the Gaussian.
  The two-dimensional Gaussian first moment is
  $\int z e^{-s|z|^2-i z\cdot\xi}dz=-i\pi\xi(2s^2)^{-1}e^{-|\xi|^2/(4s)}$.
  Substitution $u=|\xi|^2/(4s)$ gives $A_\alpha$. Handle $\xi=0$ separately.
  Gamma positivity gives $c_{\rm std}>0$; use the already proved positivity of \texttt{calpha} for the ratio.
\end{proof}

\begin{lemma}[Identification with the canonical two-dimensional operator]
  \label{lem:c13_fourier_identification}
  \lean{GSQG.c13_fourier_identification}
  \uses{def:c13_kernel_data, def:c13_normalization}
  For every real Schwartz function $f$ on $E$, the integral
  $B_{\alpha,c}f(X)=\int K(z)(f(X-z)-f(X))dz$ converges absolutely, belongs to $L^1(E;E)$,
  and has Fourier multiplier $-icA_\alpha J\xi|\xi|^{\alpha-2}$, componentwise.
  Equivalently,
  $B_{\alpha,c}f=-cA_\alpha\nabla^\perp(-\Delta)^{-1+\alpha/2}f$.
  For bounded densities the adjoint pairing and renormalized potential give the same distributional extension.
\end{lemma}
\begin{proof}
  \uses{lem:c13_test_translation, lem:c13_adjoint_integrable, lem:c13_kernel_symbol, lem:c13_renormalized_potential}
  Fourier-transform the absolutely integrable difference kernel using Fubini, or insert Schwartz
  inversion and the symbol integral. The bound $C|\xi|^{\alpha-1}$ justifies the latter exchange.
  Differentiation of the inverse transform of $|\xi|^{\alpha-2}\widehat f$ is legitimate by absolute
  integrability. Duality is meaningful for bounded densities because $B^*\Phi\in L^1$;
  no divergent unsubtracted potential or undefined product of tempered distributions is used.
\end{proof}

\section{Graph geometry and finite-domain calculus}

\begin{definition}[Admissible graphs, patches and canonical representatives]
  \label{def:c13_graph_data}
  \lean{GSQG.c13_graph_class, GSQG.c13_patch, GSQG.c13_graph_velocity, GSQG.c13_normal_speed}
  \uses{def:c13_kernel_data}
  Let $I$ be a nonempty open time interval. A graph $F:\R\times I\to\R$ is admissible if
  it is jointly smooth on $\R\times I$ (no smoothness is required at the endpoints of $I$ or outside $\R\times I$)
  and for every compact interval $J_0\subset I$ there is $M_{J_0}\ge0$
  with $|F_x(x,t)|\le M_{J_0}$ for all $x\in\R,t\in J_0$.
  No global bound on $F$, $F_{xx}$ or time derivatives at spatial infinity is assumed.
  Put $\gamma_t(s)=(s,F(s,t))$, $\tau_t(s)=(1,F_x(s,t))$,
  $\Omega_t=\{(x,y):y<F(x,t)\}$, and $\theta_F=\mathbf1_{\Omega_t}$.
  Define, for $X\notin\Gamma_t=\gamma_t(\R)$,
  \[
    U_{\alpha,c}[F](X,t)=-c\int_\R\tau_t(s)|X-\gamma_t(s)|^{-\alpha}ds,
  \]
  and give $U$ value zero on $\Gamma_t$. This candidate boundary representation is not assumed
  to satisfy the two-dimensional velocity law; that is proved in \cref{sec:c13_velocity_law}.
  Define separately
  \[
    W_{\alpha,c}[F](x,t)=c\int_\R
       [F_x(x,t)-F_x(s,t)]|\gamma_t(x)-\gamma_t(s)|^{-\alpha}ds.
  \]
  The graph normal is $(-F_x,1)/\sqrt{1+F_x^2}$; $W$ is its normal velocity multiplied by
  $\sqrt{1+F_x^2}$. It is not the value of the arbitrarily extended vector field on the graph.
\end{definition}

\begin{lemma}[Measurability, null graphs and a distance bound]
  \label{lem:c13_graph_geometry}
  \lean{GSQG.c13_graph_geometry}
  \uses{def:c13_graph_data}
  Each patch is measurable, each graph has two-dimensional measure zero, and the spacetime graph
  has three-dimensional measure zero. For a fixed-time slope bound $|F_x|\le M$, let
  $h=s-x$ and $\delta=y-F(x,t)$. Then
  \[
    h^2+\delta^2\le2(1+M^2)\bigl(h^2+(y-F(s,t))^2\bigr).
  \]
  The same constants work on a compact time interval with its common slope bound.
\end{lemma}
\begin{proof}
  Openness of the subgraph follows from continuity; Fubini makes each vertical graph section a
  null singleton. Use $|F(s,t)-F(x,t)|\le M|h|$ and the squared triangle inequality.
\end{proof}

\begin{lemma}[Finite graph-domain flux formula]
  \label{lem:c13_finite_graph_flux}
  \lean{GSQG.c13_finite_graph_flux}
  \uses{def:c13_kernel_data}
  Let $R>0$, $F\in C^1(\R)$, $H>0$ with $-H<F(s)$ for $|s|\le R$, and
  $Q:E\to E$ be $C^1$ on a neighborhood of the closed domain
  $D=\{(s,y):|s|<R,-H<y<F(s)\}$. Then
  \[
  \begin{split}
    \int_D\operatorname{div}Q
      ={}&\int_{-R}^{R}(Q_2-F'Q_1)(s,F(s))ds-\int_{-R}^{R}Q_2(s,-H)ds\\
       &+\int_{-H}^{F(R)}Q_1(R,y)dy-\int_{-H}^{F(-R)}Q_1(-R,y)dy .
  \end{split}
  \]
  Every integral is absolutely convergent.
\end{lemma}
\begin{proof}
  Coordinate Fubini and vertical FTC give the $Q_2$ terms.
  Differentiate $s\mapsto\int_{-H}^{F(s)}Q_1(s,y)dy$, then use one-dimensional FTC.
  This is the graph-specific substitute for a general Gauss theorem.
  Lean route: for the $x$-derivative term substitute $y=(F(s)+H)\tau-H$, which maps the domain to the
  rectangle $[-R,R]\times[0,1]$, then use Fubini there and FTC in $s$ and in $\tau$.
\end{proof}

\begin{lemma}[Flux formula for an unbounded subgraph and a compact test field]
  \label{lem:c13_compact_graph_flux}
  \lean{GSQG.c13_compact_graph_flux}
  \uses{def:c13_kernel_data}
  For a $C^1$ graph with bounded slope and a compactly supported $C^1$ vector field $Q$,
  \[
    \int_{y<F(x)}\operatorname{div}Q(X)dX
      =\int_\R(Q_2-F'Q_1)(x,F(x))dx,
  \]
  with absolute convergence.
\end{lemma}
\begin{proof}
  \uses{lem:c13_finite_graph_flux}
  Use $H_R=|F(0)|+(M+1)R+1$. Once $R$ is large enough, the artificial sides lie outside
  the support of $Q$. Pass to the limit in the compactly supported volume integral.
\end{proof}

\section{The candidate velocity as an actual function}

\begin{lemma}[Off-graph velocity convergence and size]
  \label{lem:c13_velocity_bound}
  \lean{GSQG.c13_velocity_bound}
  \uses{def:c13_graph_data}
  For $|F_x|\le M$ the contour integral for $U(X,t)$ converges absolutely at every $X\notin\Gamma_t$.
  There is $C_{\alpha,c,M}>0$, independent of $F,X$, such that
  \[
    |U((x,y),t)|\le C_{\alpha,c,M}|y-F(x,t)|^{1-\alpha}\quad(y\ne F(x,t)).
  \]
\end{lemma}
\begin{proof}
  Bound $|\tau|$ by $\sqrt{1+M^2}$ and compare the denominator with
  $h^2+\delta^2$, using $|F(s,t)-F(x,t)|\le M|s-x|$ directly. The one-dimensional integral of $(h^2+\delta^2)^{-\alpha/2}$
  scales as $|\delta|^{1-\alpha}$ and converges because $\alpha>1$.
\end{proof}

\begin{lemma}[Local integrability and joint contour-test Fubini]
  \label{lem:c13_velocity_local_integrable}
  \lean{GSQG.c13_velocity_local_integrable}
  \uses{def:c13_graph_data}
  For an admissible graph, $U$ is jointly measurable and locally integrable on $E\times I$.
  For a compact spatial set $Q$ and compact time interval $J_0\subset I$,
  $\int_{J_0}\int_Q|U|<\infty$.
  For every compactly supported spatial test $\Phi$, the integrand
  $|\tau_t(s)||X-\gamma_t(s)|^{-\alpha}|\Phi(X)|$ is integrable in $(s,X)$;
  its bound is uniform in $t\in J_0$ for tests supported in a fixed compact set with bounded sup norm.
\end{lemma}
\begin{proof}
  \uses{lem:c13_velocity_bound}
  Integrate $|y-F(x,t)|^{1-\alpha}$ over bounded vertical intervals; the bound is uniform
  in the center $F(x,t)$ since $1-\alpha>-1$. For joint contour-test integrability, use Tonelli:
  bounded $s$ uses local two-dimensional integrability of $|X-\gamma(s)|^{-\alpha}$;
  large $|s|$ uses horizontal separation and the integrable tail $|s|^{-\alpha}$.
  Lean route: the vertical estimate and the graph nullness are proved directly in coordinates
  (Tonelli in \texttt{lintegral}); joint measurability on $E\times I$ uses a piecewise modification,
  since $F$ is not controlled for $t\notin I$.
  Do not replace this Tonelli argument by a pointwise integral assigned zero on the singular graph.
\end{proof}

\begin{lemma}[Spatial smoothness and divergence-free velocity]
  \label{lem:c13_velocity_divergence}
  \lean{GSQG.c13_velocity_divergence}
  \uses{def:c13_graph_data}
  For each fixed time, $U$ is spatially $C^1$ on $E\setminus\Gamma_t$ and
  $\operatorname{div}_X U=0$ there. Its first spatial derivatives are jointly continuous in $(X,t)$
  off the spacetime graph. No time differentiability of $U$ is claimed from the present hypotheses.
\end{lemma}
\begin{proof}
  \uses{lem:c13_velocity_bound}
  On off-graph compact subsets the vertical gap is bounded away from zero and the global slope
  is bounded uniformly in time. The kernel and its first spatial derivative have tails
  bounded by $C|s|^{-\alpha}$ and $C|s|^{-\alpha-1}$, respectively. For divergence, use
  $\operatorname{div}U=c\int\partial_s|X-\gamma_t(s)|^{-\alpha}ds$.
  FTC on $[-R,R]$ and horizontal separation make both endpoint terms vanish.
\end{proof}

\section{Identification of the candidate with the two-dimensional velocity law}
\label{sec:c13_velocity_law}

\begin{lemma}[Artificial-boundary errors vanish]
  \label{lem:c13_truncation_boundary_error}
  \lean{GSQG.c13_truncation_boundary_error}
  \uses{def:c13_graph_data}
  Fix a time and $|F'|\le M$. Put $H_R=|F(0)|+(M+1)R+1$ and
  $D_R=\{|x|<R,-H_R<y<F(x)\}$.
  For every vector test $\Phi$, the sum of the three artificial-boundary contributions
  in the finite graph flux formula for $Q=J A_\Phi$ has modulus at most
  $C_{\alpha,c,M,F(0),\Phi}R^{1-\alpha}$ for all sufficiently large $R$.
  The constant and lower threshold for $R$ are chosen independently of $R$.
\end{lemma}
\begin{proof}
  Each artificial side is at distance at least $R/2$ from the test support once $R$ is large.
  Its length is at most $C_{M,F(0)}R$. Multiply this length by the $R^{-\alpha}$ potential bound.
  Lean route: that far-field bound, $|A_\Phi(Y)|\le|c|2^\alpha|Y|^{-\alpha}\|\Phi\|_1$ for
  $|Y|\ge2L$, is proved directly from $|X-Y|\ge|Y|/2$ on the test support, without citing the earlier lemmas.
  The exponent $1-\alpha<0$ is essential here.
\end{proof}

\begin{lemma}[Bulk-to-contour velocity identity]
  \label{lem:c13_bulk_contour_identity}
  \lean{GSQG.c13_bulk_contour_identity}
  \uses{def:c13_kernel_data, def:c13_graph_data}
  For a fixed smooth bounded-slope graph, every vector test $\Phi$ satisfies
  \[
    \int_E\theta_F(Y)B^*\Phi(Y)dY
      =-\int_\R A_\Phi(\gamma(s))\cdot\tau(s)ds
      =\int_E U[F](X)\cdot\Phi(X)dX.
  \]
  All integrals are absolutely convergent. Hence the candidate contour velocity is a locally
  integrable representative of the independently defined two-dimensional operator $B_{\alpha,c}\theta_F$.
\end{lemma}
\begin{proof}
  \uses{lem:c13_test_potential, lem:c13_finite_graph_flux, lem:c13_truncation_boundary_error, lem:c13_velocity_local_integrable, lem:c13_adjoint_integrable, lem:c13_graph_geometry}
  Apply the finite flux formula to $J A_\Phi$ on $D_R$ and use $B^*\Phi=-\operatorname{div}(JA_\Phi)$.
  On the graph, $n\,ds=(-F',1)ds=J\tau\,ds$, so $(JA_\Phi)\cdot n\,ds=A_\Phi\cdot\tau\,ds$.
  The outer minus sign gives the negative contour velocity. The artificial sides vanish by the
  preceding lemma; dominated convergence applies to $\theta_{D_R}B^*\Phi$, which is dominated by
  $|B^*\Phi|\in L^1$. The graph is null, so the value of $U$ on it is irrelevant.
  Joint contour-test Fubini gives the second equality.
\end{proof}

\begin{lemma}[Bulk principal values agree with the weak velocity]
  \label{lem:c13_bulk_principal_value}
  \lean{GSQG.c13_bulk_principal_value}
  \uses{def:c13_kernel_data, def:c13_graph_data}
  Set $u_\varepsilon(X)=\int_{|z|>\varepsilon}K(z)\theta_F(X-z)dz$ for $\varepsilon>0$.
  Each integral is absolutely convergent. Off the graph, the limit as $\varepsilon\downarrow0$
  exists and equals the absolutely convergent desingularized integral
  $\int K(z)(\theta_F(X-z)-\theta_F(X))dz$.
  The functions $u_\varepsilon$ converge in $L^1_{\rm loc}(E)$ to $U[F]$; the desingularized
  expression and $U[F]$ therefore agree a.e. No boundary vector principal value is asserted.
\end{lemma}
\begin{proof}
  \uses{lem:c13_radial_majorants, lem:c13_graph_geometry, lem:c13_adjoint_integrable, lem:c13_bulk_contour_identity, lem:c13_velocity_local_integrable}
  Oddness gives $\int_{|z|>\varepsilon}K(z)dz=0$. Off the graph the density is constant in a ball
  of radius comparable to the vertical gap, so the desingularized integrand vanishes there.
  Uniformly in $\varepsilon$, its norm is bounded by $C|y-F(x)|^{1-\alpha}$, an $L^1_{\rm loc}$
  majorant: it is bounded by a multiple of the contour majorant $\int|\tau||X-\gamma(s)|^{-\alpha}ds$,
  which is locally integrable by joint contour-test Fubini. The limit represents $B\theta_F$ by the adjoint pairing and thus equals $U$ a.e.
  Use the locally-integrable test-function uniqueness theorem, not an unproved pointwise trace claim.
\end{proof}

\section{The finite normal trace}

\begin{lemma}[Normal-speed convergence and continuity]
  \label{lem:c13_normal_speed_integrable}
  \lean{GSQG.c13_normal_speed_integrable}
  \uses{def:c13_graph_data}
  The integral defining $W[F](x,t)$ converges absolutely and $W$ is continuous on $\R\times I$.
  For $|x|\le R$, $t\in J_0\Subset I$, a bound for the integrand in the variable $h=s-x$ is
  \[
    c L_{R,J_0}|h|^{1-\alpha}\mathbf1_{0<|h|\le1}
      +2c M_{J_0}|h|^{-\alpha}\mathbf1_{|h|>1},
  \]
  where $L_{R,J_0}$ bounds $|F_{xx}|$ on $[-R-1,R+1]\times J_0$.
\end{lemma}
\begin{proof}
  Use the mean-value bound for the numerator near $s=x$ and the global slope bound in the tail.
  The denominator is at least $|s-x|^\alpha$. Shift to $h=s-x$ before dominated convergence.
\end{proof}

\begin{lemma}[Two-sided vertical normal trace with an error estimate]
  \label{lem:c13_normal_trace}
  \lean{GSQG.c13_normal_trace}
  \uses{def:c13_graph_data}
  For $d\ne0$ put
  $W_d(x,t)=U_2(x,F(x,t)+d,t)-F_x(x,t)U_1(x,F(x,t)+d,t)$.
  For every $R>0$, $J_0\Subset I$ there is $C=C(\alpha,c,R,J_0,F)>0$ such that
  \[
    \sup_{|x|\le R,\,t\in J_0}|W_d(x,t)-W(x,t)|\le C|d|^{2-\alpha}
       \quad(0<|d|\le1).
  \]
  Thus both vertical normal traces exist and agree with $W$.
\end{lemma}
\begin{proof}
  \uses{lem:c13_velocity_bound}
  Combine the two velocity components inside their absolutely convergent off-graph integrals.
  The resulting numerator is $c(F_x(x,t)-F_x(x+h,t))$.
  Split $|h|$ into $(0,|d|)$, $(|d|,1)$ and $(1,\infty)$.
  Use the near majorant in the first part and the kernel derivative bound
  $|\partial_d(h^2+(F(x,t)+d-F(x+h,t))^2)^{-\alpha/2}|\le C_\alpha|h|^{-\alpha-1}$
  in the other two. The bounds involve only the same compact-time slope and local curvature constants.
  Lean route: the convergence of the limit integral $W$ is re-derived from the same pointwise majorant,
  after the substitution $h=|d|r$, so the separate continuity lemma for $W$ is not needed here.
\end{proof}

\begin{lemma}[Flat-graph sign and trace check]
  \label{lem:c13_flat_graph_check}
  \lean{GSQG.c13_flat_graph_check}
  \uses{def:c13_graph_data}
  For $F=0$ and $d\ne0$,
  \[
    U(x,d)=-c\,C_\alpha|d|^{1-\alpha}(1,0),\qquad
    C_\alpha=\int_\R(1+s^2)^{-\alpha/2}ds\in(0,\infty),\qquad W=0.
  \]
  In particular the tangential velocity diverges in magnitude at the boundary while its normal
  component is zero. Assigning any values to $U$ on the null graph leaves its weak equations unchanged.
\end{lemma}
\begin{proof}
  Substitute $s=x+|d|r$ in the velocity integral. Positivity and tail integrability give $C_\alpha>0$.
  Use nullness of the graph to justify representative independence.
\end{proof}

\section{Transport and equivalence with the front equation}

\begin{definition}[Weak patch system with coefficient $c$]
  \label{def:c13_weak_patch_system}
  \lean{GSQG.c13_scalar_test, GSQG.c13_vector_test, GSQG.c13_velocity_law, GSQG.c13_weak_patch_system}
  \uses{def:c13_kernel_data, def:c13_graph_data}
  A scalar (respectively vector) spacetime test on $E\times I$ is a smooth compactly supported function
  $E\times\R\to\R$ (respectively $E\times\R\to E$) whose support lies in $E\times I$.
  A bounded measurable density $\theta:E\times I\to\R$ and a measurable
  $u\in L^1_{\rm loc}(E\times I;E)$ obey the coefficient-$c$ velocity law if, for every
  vector spacetime test $\Phi$,
  \[
    \int_I\!\int_E u(X,t)\cdot\Phi(X,t)dXdt
      =\int_I\!\int_E\theta(Y,t)B^*_{\alpha,c}(\Phi(\cdot,t))(Y)dYdt.
  \]
  They solve the weak patch system if also, for every scalar spacetime test $\zeta$,
  \[
    \int_I\!\int_E\theta(X,t)
       [\partial_t\zeta(X,t)+u(X,t)\cdot\nabla_X\zeta(X,t)]dXdt=0.
  \]
  All pairings must be absolutely integrable. For $\theta=\theta_F$ boundedness is automatic.
  With $c=c_{\rm std}$ this is the canonical gSQG operator, as identified in \cref{sec:c13_operator}.
  The product $\theta u$ is formed only after obtaining the locally integrable velocity representative.
\end{definition}

\begin{lemma}[Spacetime velocity law and uniqueness of the representative]
  \label{lem:c13_velocity_law_unique}
  \lean{GSQG.c13_velocity_law_unique}
  \uses{def:c13_weak_patch_system}
  For every admissible $F$, the representative $U[F]$ satisfies the coefficient-$c$ velocity law
  on $E\times I$. Every other locally integrable velocity satisfying that law for $\theta_F$
  equals $U[F]$ a.e. in spacetime. Such velocities are divergence-free distributionally in space.
\end{lemma}
\begin{proof}
  \uses{lem:c13_bulk_contour_identity, lem:c13_velocity_local_integrable, lem:c13_adjoint_integrable, lem:c13_test_potential}
  Integrate the spatial identity in time; compact-time slope bounds and test-function seminorms
  justify Fubini. Apply locally-integrable test-function uniqueness componentwise to the difference.
  For divergence use gradient tests and $B^*(\nabla\zeta)=0$.
\end{proof}

\begin{lemma}[Spatial weak flux across the actual graph]
  \label{lem:c13_spatial_flux}
  \lean{GSQG.c13_spatial_flux}
  \uses{def:c13_graph_data}
  For every fixed $t\in I$ and scalar spatial test $\zeta$,
  \[
    \int_{\Omega_t}U(X,t)\cdot\nabla\zeta(X)dX
      =\int_\R W(x,t)\zeta(x,F(x,t))dx.
  \]
  Both integrals are absolutely convergent, with estimates uniform on compact time intervals
  for tests with a common compact support and bounded first derivatives.
\end{lemma}
\begin{proof}
  \uses{lem:c13_compact_graph_flux, lem:c13_velocity_divergence, lem:c13_velocity_local_integrable, lem:c13_normal_trace, lem:c13_normal_speed_integrable}
  First integrate over $y<F(x,t)-d$, $d>0$. To apply the compact-field formula, multiply
  $U\zeta$ by a smooth function of the vertical gap that equals one for gap at least $d$
  and zero for gap at most $d/2$. This field is $C^1$ across the excluded original graph.
  On the truncated domain its divergence is $U\cdot\nabla\zeta$.
  Pass $d\downarrow0$: local integrability controls the volume term, and the cancelled normal
  trace controls the boundary term, with continuity of $W$ for the dominated boundary limit.
  Never pass to separate vector boundary traces.
\end{proof}

\begin{lemma}[Time derivative of the moving-patch pairing]
  \label{lem:c13_moving_patch_pairing}
  \lean{GSQG.c13_moving_patch_pairing}
  \uses{def:c13_graph_data}
  For every scalar spacetime test $\zeta$,
  \[
    \int_I\!\int_E\theta_F\partial_t\zeta
       =-\int_I\!\int_\R F_t(x,t)\zeta(x,F(x,t),t)dxdt,
  \]
  with absolute convergence.
\end{lemma}
\begin{proof}
  Write the spatial pairing as $\int dx\int_{-H}^{F(x,t)}\zeta(x,y,t)dy$, with $H$ below
  the vertical test support and below $F$ on its compact horizontal/time projection.
  Differentiate by endpoint FTC and parameter differentiation. Integrate the resulting total
  time derivative over the compact test support; its endpoint terms vanish.
  Lean route: translate $y=F(x,t)+s$ and set $Z(x,s,t)=\zeta(x,F(x,t)+s,t)$, so that
  $\partial_t\zeta(x,F+s,t)=\partial_tZ-F_t\partial_sZ$; then Fubini and FTC on $s\le0$ and in $t$,
  with no differentiation under the integral sign.
\end{proof}

\begin{lemma}[Lifting tests from the graph]
  \label{lem:c13_trace_tests}
  \lean{GSQG.c13_trace_tests}
  \uses{def:c13_graph_data}
  For every scalar test $\psi\in C_c^\infty(\R\times I)$ there is a scalar spacetime test
  $\zeta\in C_c^\infty(E\times I)$ satisfying $\zeta(x,F(x,t),t)=\psi(x,t)$.
  Consequently a continuous $r:\R\times I\to\R$ whose integral against all these traces is zero
  must vanish pointwise.
\end{lemma}
\begin{proof}
  Choose a fixed smooth compactly supported $\eta$ with $\eta(0)=1$ and set
  $\zeta(x,y,t)=\psi(x,t)\eta(y-F(x,t))$. The function $F$ is bounded on the compact
  support of $\psi$, making the lifted support compact. Its time support stays inside $I$,
  so extension by zero outside $I$ is smooth. Finish by test-function uniqueness
  and continuity on the open domain.
\end{proof}

\begin{theorem}[Two-dimensional patch equation is equivalent to the graph equation]
  \label{thm:c13_patch_front_equivalence}
  \lean{GSQG.c13_patch_front_equivalence}
  \uses{def:c13_weak_patch_system}
  Let $F$ be admissible. Then $(\theta_F,U_{\alpha,c}[F])$ solves the coefficient-$c$ weak patch
  system if and only if
  \[
    F_t(x,t)=c\int_\R
      \frac{F_x(x,t)-F_x(s,t)}{((x-s)^2+(F(x,t)-F(s,t))^2)^{\alpha/2}}ds
      \quad(x\in\R,t\in I).
  \]
  The same equivalence holds with any locally integrable $u$ satisfying the velocity law,
  interpreted up to a.e. equality. Every front integral is absolutely convergent.
\end{theorem}
\begin{proof}
  \uses{lem:c13_velocity_law_unique, lem:c13_spatial_flux, lem:c13_moving_patch_pairing, lem:c13_trace_tests, lem:c13_normal_speed_integrable}
  The transport pairing equals $\int(W-F_t)\zeta(x,F,t)dxdt$.
  One direction follows by substitution; in the other use lifted tests and continuity of $W-F_t$.
  A.e. equality of velocities preserves all transport pairings because $\theta_F$ is bounded.
\end{proof}

\begin{lemma}[Local time continuity and the initial-value formulation]
  \label{lem:c13_patch_time_trace}
  \lean{GSQG.c13_patch_time_trace}
  \uses{def:c13_weak_patch_system}
  For $R>0$ and $s,t\in I$,
  \[
    \int_{[-R,R]^2}|\theta_F(X,t)-\theta_F(X,s)|dX
       \le\int_{-R}^R|F(x,t)-F(x,s)|dx.
  \]
  Thus $t\mapsto\theta_F(t)$ is continuous into $L^1_{\rm loc}(E)$.
  If the equivalent equations above hold and $t_0<t_1$ both lie in $I$, let $\zeta$ be smooth
  on a neighborhood of $E\times[t_0,t_1]$, have a common compact spatial support there,
  and vanish for $t$ sufficiently close to $t_1$. It may be nonzero at $t_0$.
  Then
  \[
    \int_{t_0}^{t_1}\!\int_E\theta_F(\partial_t\zeta+U\cdot\nabla\zeta)
       +\int_E\theta_F(X,t_0)\zeta(X,t_0)dX=0.
  \]
  For an initial endpoint outside $I$, this last assertion requires additionally a locally uniform
  graph limit and a common global slope bound on a one-sided time neighborhood; it is not implied
  by interior compact-time bounds alone.
\end{lemma}
\begin{proof}
  \uses{lem:c13_velocity_local_integrable, lem:c13_velocity_bound}
  Integrate the difference of two half-line indicators on vertical slices.
  Uniform continuity of $F$ on compact horizontal/time sets gives the local time trace.
  For the initial-value identity use only the weak transport equation, tested with
  $\chi_h(t)\zeta$ for a time cutoff $\chi_h$ rising from $0$ at $t_0$ to $1$ at $t_0+h$, together with the
  $L^1_{\rm loc}$ continuity just proved.
  Retain the lower endpoint term in the time pairing identity; its sign is positive in the weak formula.
  Initial-value tests are restrictions of smooth ambient functions, not globally smooth functions
  required to vanish before $t_0$; that latter restriction would force the initial term to be zero.
  For the stated endpoint extension, start at $t_0+h$ and use the uniform-slope integrability
  estimate and the $L^1_{\rm loc}$ density trace as $h\downarrow0$.
  Lean route: for that endpoint case, the velocity bound and the common slope bound give integrability
  of $U$ on $K\times(t_0,t_1]$; the limit $F_0$ is continuous, and the $L^1_{\rm loc}$ density limit comes
  from the slice estimate above applied to the glued graph equal to $F_0$ at $t_0$ and to $F$ for $t>t_0$.
\end{proof}

\section{Transfer of the front theorem to the canonical patch system}

\begin{lemma}[Uniform slope for the self-similar family]
  \label{lem:c13_main_graph_admissible}
  \lean{GSQG.c13_main_graph_admissible}
  \uses{def:c13_graph_data, def:self_similar_ansatz}
  For every profile supplied by \cref{thm:main}, the graph $F_{\epsilon,T}$ is admissible
  on $(-\infty,T)$ and satisfies
  $|F_x(x,t)|\le\sqrt\epsilon(C+1)$ for all $x,t<T$, with the profile-bound constant $C$.
  No global curvature bound is required.
\end{lemma}
\begin{proof}
  \uses{thm:main, lem:c01_ansatz_time_derivative}
  The two profile estimates imply $|f'|\le C+1$. Differentiate the ansatz to obtain
  $F_x=\sqrt\epsilon f'(x/L(t))$. Joint smoothness is already part of the main theorem.
  Local curvature bounds follow on spacetime compacts lying strictly below $T$.
\end{proof}

\begin{lemma}[Exact time normalization]
  \label{lem:c13_time_normalization}
  \lean{GSQG.c13_time_normalization}
  \uses{def:c13_normalization, def:gsqg_front_equation, def:c13_graph_data}
  Let $p=\alpha-1$, $\kappa=\kappa_\alpha>0$ and let an admissible $F$ solve the
  normalized front equation on $I$. Then
  $G(x,t)=F(x,t/\kappa)$, on $\{t:t/\kappa\in I\}$, is admissible and satisfies the
  coefficient-$c_{\rm std}$ graph equation. If $I=(-\infty,T/\kappa)$, the new interval is
  $(-\infty,T)$. For the self-similar ansatz its physical scale is
  \[
    L_{\rm std}(t)=(\alpha(T-t)/\kappa)^{1/\alpha}.
  \]
\end{lemma}
\begin{proof}
  \uses{lem:c01_paper_front_equation}
  The literal normalized graph equation has coefficient $c_{\rm front}$.
  Time differentiation contributes $1/\kappa=c_{\rm std}/c_{\rm front}$.
  The ansatz with normalized blow-up time $T/\kappa$ gives the displayed scale.
  Positivity of $\kappa$ follows directly from positivity of the Gamma factors in $A_\alpha$ and of \texttt{calpha}.
  Preserve the density jump equal to one; do not absorb the coefficient by silently changing the patch strength.
\end{proof}

\begin{theorem}[Canonical two-dimensional patches forming a corner]
  \label{thm:c13_patch_main}
  \lean{GSQG.c13_patch_main}
  \uses{def:c13_weak_patch_system, def:c13_normalization, def:self_similar_ansatz}
  Let $\alpha\in(7/4,2)$ and $p=\alpha-1$. There exist $\epsilon_0,C>0$, depending only on
  $\alpha$, such that for every $0<\epsilon<\epsilon_0$ there is a profile $f_\epsilon$
  satisfying all profile estimates of \cref{thm:main} with constant $C$ and the following property.
  For every $T>0$, define
  \[
    G(x,t)=\sqrt\epsilon L_{\rm std}(t)f_\epsilon(x/L_{\rm std}(t)),\qquad
    \theta(X,t)=\mathbf1_{X_2<G(X_1,t)},\qquad u=U_{\alpha,c_{\rm std}}[G].
  \]
  Then on $E\times(-\infty,T)$ the pair is a weak solution of
  \[
    \partial_t\theta+\operatorname{div}(\theta u)=0,\qquad
    u=-\nabla^\perp(-\Delta)^{-1+\alpha/2}\theta,
  \]
  with the canonical normalization defined above and with the initial trace at $t=0$.
  The interface is smooth for $t<T$ and
  $\|G(\cdot,t)-\sqrt\epsilon x_-\|_\infty\le C\sqrt\epsilon(T-t)^{1/\alpha}$.
  For each such $\epsilon$ there is $b_\epsilon>0$, independent of $T$ and $\delta$, such that
  for every $T,\delta>0$ there is $t_\delta<T$ with
  $\sup_{|x|<\delta}|\partial_x^2G(x,t)|\ge b_\epsilon(T-t)^{-1/\alpha}$
  whenever $t_\delta<t<T$.
  Moreover, with $\theta_*(X)=\mathbf1_{X_2<\sqrt\epsilon(X_1)_-}$,
  $\theta(t)\to\theta_*$ in $L^1_{\rm loc}(E)$ as $t\uparrow T$.
  This is an infinite-area patch result; no global $L^1$ assumption on $\theta$ is made.
\end{theorem}
\begin{proof}
  \uses{thm:main, lem:c13_main_graph_admissible, lem:c13_time_normalization, thm:c13_patch_front_equivalence, lem:c13_fourier_identification, lem:c13_patch_time_trace}
  Apply the completed front theorem at normalized time $T/\kappa$ and use the equivalence theorem
  with coefficient $c_{\rm std}$. The initial time zero is interior to the time interval.
  Corner convergence acquires the fixed factor $\kappa^{-1/\alpha}$ and the curvature lower bound
  the fixed factor $\kappa^{1/\alpha}$; positivity preserves both claims.
  On $[-R,R]^2$, the density error is at most $2R$ times the uniform graph-height error.
  No classical vector velocity at the interface or continuation beyond the corner time is claimed.
\end{proof}
